% ==============================================================================
% Section 3: Related Works & Methodological Benchmarking
% ==============================================================================

\section{Related Works and Methodological Benchmarking}
\label{sec:related_works}

Computational models of ecosystem dynamics, plant physiology, and consumer--resource interactions have historically developed within separate scale regimes \cite{okubo2001diffusion, swat2012multiscale}. Existing platforms generally target either regional biogeochemical crop balances, coarse-grained landscape wildlife populations, or sub-cellular tissue morphogenesis. Table~\ref{tab:simulator_benchmarking} positions PHIDS relative to these frameworks across spatial resolution, PDE transport capability, and execution throughput.

\subsection{Continuous Macro Biogeochemical Models}
\label{sec:related_biogeochem}

At landscape and regional scales ($> 100\,\text{m}$), computational agronomy and forestry rely on biogeochemical suites such as SIMPLACE \cite{enders2023simplace}, MONICA \cite{nendel2011monica}, and LANDIS-II \cite{scheller2007landis}. These tools forecast crop yields, soil water transport, and carbon--nitrogen turnover using one-dimensional vertical soil columns solving Richards water percolation and Century/RothC-style mineralisation ODEs. Although well validated for abiotic crop--soil budgets, they operate without autonomous foraging agents (reducing herbivory to exogenous defoliation scalars), resolve only daily $\Delta t = 24\,\text{h}$ steps too coarse for $1\text{--}3\,\text{h}$ plume dispersion and taxis, and exclude horizontal 2D reaction--diffusion--advection PDE solvers.

\subsection{Agent-Based Landscape and Trophic Frameworks}
\label{sec:related_abm}

Spatially explicit Agent-Based Models (ABMs) and Individual-Based Models (IBMs) capture individual movement over heterogeneous habitat maps, exemplified by ALMaSS (Animal, Landscape and Man Simulation System) \cite{topping2003almass}. ALMaSS simulates wildlife demography and pesticide exposure via procedural state machines over GIS rasters. At densities of $10^4\text{--}10^5$ interacting plants and insect swarms, however, traditional object-oriented ABMs incur pointer chasing and cache misses across scattered heap instances, and ALMaSS relies on static or piecewise-constant environmental rasters without coupled continuous semiochemical PDE solvers. General-purpose frameworks such as NetLogo \cite{wilensky1999netlogo} and Mesa \cite{kazil2020mesa} are treated separately from Table~\ref{tab:simulator_benchmarking} because they are unspecialised programming toolkits lacking built-in biophysical PDE solvers, calibrated kinetic modules, or contiguous SoA memory layouts.

\subsection{Cellular Biophysics Simulators}
\label{sec:related_cpm}

At microscopic tissue scales ($1\text{--}100\,\mu\text{m}$), hybrid continuous--discrete simulators such as CompuCell3D (CC3D) \cite{swat2012multiscale}, the Hybrid Automata Library (HAL) \cite{bravo2020hybrid}, and MCell4 \cite{husar2024mcell4} couple lattice automata or molecular Brownian particles with continuum reaction--diffusion PDEs for morphogenesis and tumour biology. In CC3D, cells span many lattice pixels evolved via Monte Carlo Step spin flips under a global Hamiltonian:
\begin{equation}
  \mathcal{H} = \sum_{\text{interfaces}} J(\tau(\sigma_i), \tau(\sigma_j)) + \sum_{\text{cells}} \lambda_{\text{vol}} (v_\sigma - V_t)^2,
  \label{eq:cc3d_hamiltonian}
\end{equation}
where $J(\tau(\sigma_i), \tau(\sigma_j))$ is interfacial contact energy, $\lambda_{\text{vol}}$ volume elasticity, $v_\sigma$ cell volume, and $V_t$ target volume. Because computational effort is dominated by sub-cellular boundary fluctuations or molecular collisions, these frameworks do not scale directly to meso-scale ecosystems requiring foraging momentum, multi-rate diurnal/weekly strides, and dual airborne--mycorrhizal signalling.

\begin{table*}[b]
  \centering
  \caption{Methodological Comparison Across Domain-Specific Ecological and Biophysical Simulation Architectures.}
  \label{tab:simulator_benchmarking}
  % ==============================================================================
% Table 4: Comprehensive Methodological Comparison Across Simulators
% Placed at [b] (bottom of page)
% Narrative Story:
% This table systematically justifies the architectural necessity of PHIDS against
% established ecological, agronomic, and biophysical simulation platforms.
% It illustrates the historical compromise: existing tools either model regional
% crop yields without biotic predators (SIMPLACE, MONICA), landscape succession
% without reactive semiochemical transport (LANDIS-II, ALMaSS), or micro-scale
% tissue morphogenesis without macroscopic animals (CompuCell3D).
% Blank-canvas toolkits (NetLogo, Mesa) are excluded from the comparison matrix
% because their unspecialized nature produces uninformative "User-Defined" entries.
% ==============================================================================

\scriptsize
\renewcommand{\arraystretch}{1.18}
\begin{tabularx}{\textwidth}{
  >{\raggedright\arraybackslash\hsize=1.05\hsize}X
  >{\raggedright\arraybackslash\hsize=1.15\hsize}X
  >{\raggedright\arraybackslash\hsize=0.90\hsize}X
  >{\raggedright\arraybackslash\hsize=0.95\hsize}X
  >{\raggedright\arraybackslash\hsize=0.95\hsize}X
  >{\raggedright\arraybackslash\hsize=0.98\hsize}X
  >{\raggedright\arraybackslash\hsize=1.02\hsize}X
}
  \toprule
  \textbf{Capability / Dimension} &
  \textbf{PHIDS (This Work)} &
  \textbf{SIMPLACE} &
  \textbf{MONICA} &
  \textbf{LANDIS-II} &
  \textbf{ALMaSS} &
  \textbf{CompuCell3D} \\
  \midrule
  \textbf{Primary Domain} &
  Meso-scale plant--herbivore semiochemical dynamics &
  Crop phenology \& agricultural yields &
  Agro-ecosystem carbon \& nitrogen kinetics &
  Forest succession \& disturbance ecology &
  Wildlife landscape population dynamics &
  Cellular tissue morphogenesis \& tumours \\
  \addlinespace[2.5pt]
  \textbf{Spatial Topology} &
  2D Toroidal $\mathbb{T}^2$ (power-of-two grids $32^2\text{--}1024^2$) &
  1D Vertical soil profile columns &
  1D Vertical layered soil columns &
  2D Raster grid ($10\text{--}100\,\text{m}$ cells) &
  2D Heterogeneous vector/raster landscape &
  2D/3D Cellular lattice ($\mu\text{m}$ domain) \\
  \addlinespace[2.5pt]
  \textbf{Temporal Stride} &
  Multi-rate strides ($1\,\text{h}, 24\,\text{h}, 168\,\text{h}$) &
  Daily discrete steps ($\Delta t = 24\,\text{h}$) &
  Daily agronomic steps ($\Delta t = 24\,\text{h}$) &
  Seasonal to multi-year increments &
  Daily to hourly animal steps &
  Monte Carlo Steps (MCS sub-second) \\
  \addlinespace[2.5pt]
  \textbf{Entity Architecture} &
  Data-Oriented SoA ECS (Contiguous field buffers) &
  Monolithic Java class hierarchy &
  Procedural C++ crop modules &
  Object-oriented C\# cohorts &
  Object-oriented C++ polymorphic agents &
  Cellular Potts Model (Generalised Ising) \\
  \addlinespace[2.5pt]
  \textbf{Continuous Fields} &
  Operator-split 2D Reaction-Diffusion PDEs &
  None (1D Richards ODE water flow) &
  None (1D ODE nitrogen percolation) &
  None (Coarse spatial rasters) &
  None (Static GIS habitat maps) &
  Finite difference diffusion solver \\
  \addlinespace[2.5pt]
  \textbf{Defence Signalling} &
  Airborne VOC plumes \& non-linear Hill priming &
  None (Abiotic crop growth only) &
  None (Abiotic fertiliser response) &
  None (Static species growth tables) &
  None (Static vegetation foraging map) &
  Local morphogen chemical secretion \\
  \addlinespace[2.5pt]
  \textbf{Subterranean Conduits} &
  Cardinal-adjacent dynamic root-link graph ($\mathbf{A}_{uv}$) &
  None &
  None &
  None &
  None &
  None (Contact energy matrix only) \\
  \addlinespace[2.5pt]
  \textbf{Foraging Kinematics} &
  Jacobi chemotaxis, MVT-inspired exit heuristic, heading persistence &
  None (No animal consumers) &
  None (No animal consumers) &
  Stochastic dispersal kernel &
  Rule-based behavioural state machines &
  Hamiltonian energy minimisation $\Delta \mathcal{H}$ \\
  \addlinespace[2.5pt]
  \textbf{Hardware \& Execution} &
  Compiled SIMD kernels ($1.25\text{--}2.59\,\text{ms}$/tick; $64\times 32\text{--}128^2$) &
  JVM threaded multi-run execution &
  Single-threaded native C++ executable &
  .NET runtime batch processor &
  Single-threaded native C++ engine &
  C++/OpenMP parallel Potts solver \\
  \addlinespace[2.5pt]
  \textbf{Replay Buffering} &
  Deterministic chunked Zarr array replay &
  Tabular CSV / SQL relational loggers &
  NetCDF / CSV output text dump &
  Output raster GIS files &
  Custom binary event logs &
  VTK / image frame sequential dumps \\
  \bottomrule
\end{tabularx}

\end{table*}

\subsection{PHIDS Architectural Niche and Throughput}
\label{sec:related_phids_niche}

Addressing meso-scale ($1\text{--}100\,\text{m}$) semiochemical plant--herbivore dynamics alongside abstract energy-allocation trade-offs relevant to sensor-network DSE requires combining continuum transport with high-density discrete ecology (Table~\ref{tab:simulator_benchmarking}) \cite{dietrich2025visual}. PHIDS bridges this gap through three architectural choices: (i) storing continuous fields and species slices in contiguous Structure-of-Arrays (SoA) NumPy buffers compiled via Numba with zero inner-loop heap allocation \cite{harris2020array, lam2015numba}; (ii) solving 2D parabolic VOC advection--diffusion on a double-buffered torus $\Omega = \mathbb{T}^2$ for thread-safe determinism (Section~\ref{subsec:double_buffering_zarr}); and (iii) decoupling hourly taxis ($\Delta t = 1\,\text{h}$), diurnal feeding ($24\,\text{h}$), and phase-staggered weekly plant growth ($168\,\text{h}$ cohorts), reducing per-tick growth memory write traffic by $168\times$ while preserving contiguous SIMD vectorisation. Benchmark measurements detailed in Section~\ref{sec:computational_realisation} sustain median tick latencies of $1.25\text{--}2.59\,\text{ms}$ on $64 \times 32$ and $128 \times 128$ grids ($84\text{--}92\%$ frame headroom relative to a $16.67\,\text{ms}$ $60\,\text{Hz}$ budget), $36.6\text{--}57.2\,\text{ms}$ on 3-species $256 \times 256$ canopies, and $137.0\,\text{ms}$ mean per tick on the full 11-species, 9-channel $256 \times 256$ ecotone.
